\section{Conclusion}\label{sec:conclusion}

We presented a two-stage pipeline for detecting EPBs in $630.0\text{ nm}$ all-sky images. A sky quality check based on star counts removes cloud-covered frames, and the remaining frames go through background subtraction, CLAHE, 2DPCA, and a stacking ensemble, at a cost of about $0.6\text{ ms}$ per image on a CPU.

Three points stand out from the experiments. First, a high score on random splits is not enough to show that a model has learned the phenomenon. The version without radiometric preprocessing had the best holdout accuracy and missed every EPB at the unseen station, whereas the full pipeline detected all of them there and stayed above $95\%$ accuracy when whole time windows or the corner text were removed. The ablation suggests that background normalization matters more for transfer between sites than the choice of classifier. Second, cloud screening is needed in practice, since contrast enhancement turns cloud edges into false bubbles; Stage~1 removed most of these false alarms on an overcast night while discarding few frames on clear and active nights. Third, reconstruction residuals and SHAP maps place the evidence for an EPB on the depletion bands, not on the image border or the printed text.

The study also has limitations. The dataset is small and covers a limited period; CP contributed only positive examples, so specificity could not be measured in the zero-shot test; and all bright-moon images are negatives. Future work will add negative examples from low-latitude stations such as CP, include several years of observations covering different phases of the solar cycle, and combine the optical classification with collocated GNSS $S_4$ scintillation receivers and digisondes for real-time space weather monitoring.
